\section{}

\comment{
The article is interesting and has an interesting premise but needs additional work.
I have included a number of grammatical edits below. Please make these corrections and ask a peer to review the article for other edits that I may have missed. I included non-grammatical comments below the grammatical fixes.
P2, L23: “his/her” not its
P2, L29: “to estimate”
P3, L22-23: inconsistent quotations
P4, L40: characterizes
P4, L50: capitalize Euclidean
P5, L36: rewrite sentence
L38: “short corners?” (There are only two corners and they are both short.)
41: English
42: English
43: will “shoot”
Break up the first paragraph in 4.1 into multiple smaller ones.
}

\reply{
	We thank the reviewer for their comments. We corrected the grammatical errors. Regarding the ``short corners'', the ``short corners'' refers to specific location on the court. We added a graphic that explain the different names of the court location in the Appendix.
}

\comment{
Can you explain how this is shown by the 3rd PC?
``Interestingly, players who perform well near the basket also tend to attempt and miss shots from the corners''
}

\reply{
	The explanation of the third principal components has been rewritten.
}

\comment{
Last sentence (P5, L55-56) does not make sense. What is “the axis of the hoop” or “the axis of the basket” (you use both terms)?
}

\reply{
	The explanation of the fourth principal component has also been rewritten.
}

\comment{
Figure 5: for visualization purposes, consider normalizing the scores so that the axes all line up. Also consider adding a vertical line at x = 0 so you can more easily determine if the score is positive or negative. If possible, caption each graph with the interpretation you are assigning to it so that a reader can look at the graph and understand what it means for the labeled players.
}

\reply{
	We thank the reviewer for the suggestion. We have standardized the scores and added a vertical line at $x = 0$. We also added the interpretation of the eigenfunctions that goes with the scores in the caption.
}

\comment{
P6 L51: ``designated'' as outliers
}

\reply{
	We thank the reviewer for spotting the typo.
}

\comment{
P9 L51: describe the clusters in numeric order.
}

\reply{
	The ordering of the clusters has been changed, and the clusters are now described in numeric order.
}

\comment{
P10: give a link / reference to Hawk-Eye; “player’s”
}

\reply{
The typo (players $\rigtharrow$ player's) has been fixed. The Hawkeye reference has been added.
\color{purple}\textbf{Steven}: I have provided the Bib entry for the companies web page below. Can you please cite in the appropriate place please?
}
\begin{verbatim}
    @misc{HawkeyeInnovation2025,
  title        = {Hawkeye Innovation, LLC},
  howpublished = {\url{https://www.hawkeyeinnovation.com/}},
  year         = {2025},
  organization = {Hawkeye Innovation, LLC},
  address      = {Agawam, Massachusetts, United States},
}
\end{verbatim}

\comment{
At the very least, you should report the number of guards, forwards, and centers in each cluster. (Many readers may not be able to distinguish who is a guard or who is a forward, etc, just by reading players’ names.)  It would also be nice if you used a quantitative metric to measure the agreement between the clusters you find and the NBA positions. If you try other functional clustering methods (as I suggest below), these comparisons would be even more interesting.
}

\reply{
	We added the number of guards, forwards and centers in each cluster in the appendix. We also changed the confusion matrix to show the NBA position directly. See Comment 0.4 for a discussion on metrics to assess the clustering quality.
}

\comment{
Did you consider other functional clustering approaches? It would be interesting to see if you get similar clusters (and similar levels of agreement) using a different approach.
}

\reply{See Comment 0.4 for a discussion on other functional clustering approaches.}

\comment{
Why not cluster the surfaces themselves and not the MFPCA scores using functional clustering algorithms? That would be fun to see.
}

\reply{
	See Comment 0.4 for a discussion on other functional clustering approaches.
}

\comment{
What about just FPCA of all shot locations (as opposed to made and missed shots?) Do you gain anything from doing MFPCA in the way that you do vs just FPCA of shot locations?
}

\reply{
	This is an interesting suggestion. Our rationale is that a players typical shooting ``pattern" or ``behaviour" is defined not only by where they shoot, but more specifically by where they score and miss from. In a simple example, we might consider two players who attempt a lot of shots from both the three-point line \emph{and} under the paint.
    Player $A$ is an expert three-point shooter but weaker at closer-range shots, and vice-versa for player $B$. 
    With a standard FPCA of overall shot volumes as suggested, we would characterize these player similarly. 
    However, the MFPCA of made and missed shots would would be able to decipher between the two patterns.


    As noted in Comment $0.2$, however, we do agree that future work could use additional information (e.g., overall shots volume, normalized or un-normalized) to provide improved characterisations.
}
